%% kapa: exam=1 round=21 year=2010 period=2 form=A subject=accounting subject_ko=회계학 seq=15 number=15
%% source: Q-Net 감정평가사 기출문제 (cst00309) artlSeq=5024181 file=extracted/1008/5024181/2010년 제21회 감정평가사 1차시험 기출문제/2교시.hwp
%% conversion: kordoc 4.24.51 -> tools/gen1.py (auto); review=proofread-gpt-sol; flags=none
\begin{question}{15}
(주)한국리스는 20×0년 초 (주)민국과 운용리스계약을 체결하였다. 리스계약과 관련된 주요 내용은 다음과 같다. 아래 리스계약의 회계처리가 (주)한국리스의 20×0년도 당기순이익을 얼마나 증가시키는가? (단, 리스자산의 사용효익이 감소하는 기간별 형태를 잘 나타내는 별도의 체계적 방법은 없음)
\begin{framedtext}
ㆍ(주)한국리스는 리스자산인 기계장치를 20×0년 초에 ￦200,000에 취득함과 동시에 (주)민국에 제공하였다. \par
ㆍ동 리스자산의 내용연수와 잔존가치는 각각 10년과 ￦20,000으로 추정되며, 정액법으로 감가상각한다. \par
ㆍ리스기간은 20×0년 초부터 3년간이며, 리스료는 20×0년 말에 ￦45,000, 20×1년 말에 ￦40,000, 20×2년 말에 ￦32,000을 수령하기로 하였다. \par
ㆍ(주)한국리스의 리스개설직접원가는 ￦6,000이고, 동 금액을 20×0년 초에 현금으로 지급하였다. \par
\end{framedtext}
\begin{choices}
  \item ￦15,000
  \item ￦19,000
  \item ￦21,000
  \item ￦27,000
  \item ￦30,000
\end{choices}
\end{question}
\answer{2}
